\chapter{Diseño del servicio}
\label{cap:diseno}

\section{Arquitectura general}
El servicio se compone de tres planos claramente separados: el plano de
entrada, que valida y normaliza los proyectos recibidos; el plano de
compilación, formado por contenedores efímeros; y el plano de almacenamiento,
donde se guardan los artefactos generados. La figura~\ref{fig:arquitectura}
muestra la relación entre ellos.

\begin{figure}[htbp]
  \centering
  \begin{tikzpicture}[
      font=\small, node distance=6mm and 10mm,
      caja/.style={draw, rounded corners=2pt, minimum width=28mm, minimum height=10mm,
                   align=center, fill=white, thick},
      flecha/.style={-{Latex[length=2.5mm]}, thick},
    ]
    \node[caja] (cliente) {Cliente\\(SDK o CLI)};
    \node[caja, right=of cliente, fill=principal!10] (api) {API};
    \node[caja, right=of api, fill=acento!15] (cache) {Caché de\\artefactos};
    \node[caja, below=of api, fill=green!10] (worker) {Trabajador};
    \node[caja, right=of worker, fill=red!8] (sandbox) {Contenedor\\aislado};
    \node[caja, below=of worker, fill=gray!10] (almacen) {Almacén\\de archivos};
    \draw[flecha] (cliente) -- node[above]{proyecto} (api);
    \draw[flecha] (api) -- node[above]{hash} (cache);
    \draw[flecha] (api) -- node[left]{fallo de caché} (worker);
    \draw[flecha] (worker) -- node[above]{lanza} (sandbox);
    \draw[flecha] (sandbox) |- node[pos=0.75, below]{PDF y registro} (almacen);
    \draw[flecha] (almacen) -- (worker);
    \draw[flecha, dashed] (cache) |- (almacen);
    \begin{scope}[on background layer]
      \node[fit=(worker)(sandbox)(almacen), draw=gray!60, dashed, rounded corners,
            inner sep=4pt, label={[gray]below:plano de compilación}] {};
    \end{scope}
  \end{tikzpicture}
  \caption{Arquitectura del servicio. El plano de compilación no tiene acceso a
    la red ni a los datos de otros usuarios.}
  \label{fig:arquitectura}
\end{figure}

Cada petición se procesa en cuatro pasos. Primero, la API comprueba los
límites de tamaño y rechaza rutas peligrosas (referencias al directorio
superior, rutas absolutas o enlaces simbólicos). Segundo, calcula un resumen
criptográfico del proyecto normalizado. Tercero, consulta la caché y, si
existe un artefacto para ese resumen, lo devuelve de inmediato. Cuarto, en
caso contrario, un trabajador lanza un contenedor, ejecuta la compilación y
almacena el resultado.

\section{Modelo de aislamiento}
Un documento \LaTeX{} es, en la práctica, un programa: puede leer y escribir
archivos, consumir memoria sin límite y, si el intérprete lo permite,
ejecutar órdenes del sistema. Por ello cada compilación se ejecuta en un
contenedor nuevo con las restricciones de la tabla~\ref{tab:limites}.

\begin{table}[htbp]
  \centering
  \caption{Límites aplicados a cada compilación.}
  \label{tab:limites}
  \begin{tabular}{@{}llr@{}}
    \toprule
    Recurso & Mecanismo & Límite \\
    \midrule
    Tiempo de pared & \texttt{timeout -s KILL} & \SI{60}{s} (síncrono) / \SI{300}{s} (asíncrono) \\
    Memoria & cgroup v2 & \SI{2}{GiB} \\
    Procesadores & cgroup v2 & 2 vCPU \\
    Procesos & \texttt{pids.max} & 64 \\
    Red & espacio de nombres & ninguna \\
    Llamadas al sistema & seccomp & lista blanca \\
    Ejecución de programas & \texttt{shell\_escape = f} & deshabilitada \\
    \bottomrule
  \end{tabular}
\end{table}

Además de los límites de recursos, el sistema de archivos raíz del contenedor
es de solo lectura y el directorio de trabajo se monta en memoria. La opción
\texttt{openout\_any = p} del archivo \texttt{texmf.cnf} impide que el motor
escriba fuera de ese directorio, y \texttt{-no-shell-escape} refuerza en la
línea de órdenes la prohibición ya establecida en la configuración.

\section{Clave de caché}
Sea $B = \{(p_i, b_i)\}_{i=1}^{n}$ el conjunto de archivos del proyecto, donde
$p_i$ es la ruta relativa normalizada y $b_i$ su contenido en bytes. Dado el
motor $e$ y la versión de la imagen $v$, la clave de caché se define como
\begin{equation}
  \kappa(B, e, v) = H\Big( e \,\|\, v \,\|\, \big\Vert_{i=1}^{n}
    \big(|p_i| \,\|\, p_i \,\|\, |b_i| \,\|\, b_i\big) \Big),
  \label{eq:clave}
\end{equation}
donde $H$ es SHA-256 y $\Vert$ denota la concatenación con prefijos de
longitud. Incluir $v$ en la clave garantiza que una actualización de
\TeX{} Live invalide los artefactos antiguos, que son la principal fuente de
resultados no reproducibles en proyectos de larga duración.

\section{Interfaz de programación}
El listado~\ref{lst:peticion} muestra una petición mínima al servicio. El
cuerpo contiene los archivos del proyecto codificados como texto y el nombre
del archivo principal; la respuesta es el PDF o, si se solicita mediante la
cabecera \texttt{Accept}, un objeto JSON con los diagnósticos.

\begin{lstlisting}[language=bash, caption={Petición de compilación con curl.}, label={lst:peticion}]
curl -X POST https://latex-api.example/v1/compile \
  -H "Authorization: Bearer $API_KEY" \
  -H "Content-Type: application/json" \
  -d '{
        "engine": "pdflatex",
        "entrypoint": "main.tex",
        "files": [
          { "path": "main.tex", "content": "\\documentclass{article}..." }
        ]
      }' \
  -o salida.pdf
\end{lstlisting}

Cuando la compilación falla, el servicio devuelve un documento de problema
conforme a RFC~9457 con el código \texttt{compile\_failed} y la lista de
diagnósticos, cada uno con su archivo, línea y mensaje. De este modo el
cliente puede señalar el error directamente en el editor sin analizar el
registro completo del motor.
